\scanpage{9}
où $q$ est de degré $2$ et définit l'involution $\sigma_{\widetilde A}$. Les applications $o$ et $p$ sont induites par les inclusions $\widetilde A\hookrightarrow Y$, $I\hookrightarrow Y$ en passant aux $\pi_0$.

Le système
\diagram{\begin{tikzcd}
\bigl(\widetilde A\arrow[r,"o"] & S,\sigma_{\widetilde A}\bigr)
\end{tikzcd}}
ou
\diagram{\begin{tikzcd}
\bigl(\widetilde A\arrow[r,"o"]\arrow[d,"\text{deg 2}"'] & S\bigr)\\
A &
\end{tikzcd}}
, peut être considéré comme définissant un \uline{graphe}, dont $S$ est l'ens. des sommets, et $\widetilde A$ l'ens. des \uline{arcs} (ou arêtes orientées) l'application $o$ étant l'application ``origine d'un arc''. Ce graphe ne dépend que de $\widehat X$, pas de $X$ i.e. du choix de $I\subset\widehat X(k)$. C'est pour tenir compte de ce choix que l'on considère, en plus de la structure de graphe, la donnée supplémentaire
\[
(9)\qquad I\longrightarrow S
\]
\mspage{8}
